%% ============================================================
%%  Springer Nature sn-jnl -- single column
%%  Files needed: sn-jnl.cls  sn-mathphys-num.bst
%% ============================================================
\documentclass[sn-mathphys-num]{sn-jnl}

\usepackage{graphicx}
\usepackage{amsmath,amssymb,amsfonts}
\usepackage{multirow}
\usepackage{booktabs}
\usepackage{textcomp}
\usepackage{xcolor}
\usepackage{subcaption}
\usepackage{array}
\usepackage{hyperref}
\usepackage{enumitem}
\usepackage{float}

%% Float placement
\setcounter{topnumber}{6}
\setcounter{bottomnumber}{6}
\setcounter{totalnumber}{10}
\renewcommand{\topfraction}{0.9}
\renewcommand{\bottomfraction}{0.9}
\renewcommand{\textfraction}{0.05}
\renewcommand{\floatpagefraction}{0.85}

%% ============================================================
\title[Explainable Hybrid Attention U-Net--U-Net++]{Explainable Hybrid Attention
U-Net--U-Net++ for Cross-Dataset Brain Tumor Analysis}

%% \author* marks the corresponding author; email goes in \author via \email{}
\author*[1]{\fnm{Gowripushpa} \sur{Geddam}\,\textsuperscript{a}}
\author[3]{\fnm{Jayasri} \sur{Kotti}}
\author[4]{\fnm{Bindu Madhuri} \sur{Ch.}}

\affil[1]{\orgdiv{Department of Computer Science \& Engineering},
          \orgname{JNTU-GV University},
          \orgaddress{\street{Dwarapudi}, \city{Vizianagaram},
                      \postcode{535003}, \state{Andhra Pradesh}, \country{India}}}

\affil[2]{\orgdiv{Department of Computer Science \& Engineering},
          \orgname{Anil Neerukonda Institute of Technology \& Sciences},
          \orgaddress{\city{Visakhapatnam}, \postcode{531162},
                      \state{Andhra Pradesh}, \country{India}}}

\affil[3]{\orgdiv{Department of CSE AI \& DS},
          \orgname{GMR Institute of Technology},
          \orgaddress{\city{Rajam}, \state{Andhra Pradesh}, \country{India}}}

\affil[4]{\orgdiv{Department of Information \& Technology},
          \orgname{JNTU-GV University, College of Engineering Vizianagaram},
          \orgaddress{\street{Dwarapudi}, \city{Vizianagaram},
                      \postcode{535003}, \state{Andhra Pradesh}, \country{India}}}

%% ============================================================
\begin{document}
\maketitle

%% Corresponding author footnote — placed AFTER \maketitle so it appears
%% on the title page itself, not on a blank page before it.
{\footnotesize\noindent author. E-mail:
\texttt{gowripushpa11@gmail.com}}

\begin{abstract}
Brain tumor detection from multi-modal MRI scans is demanding due to variability in
tumor shape, size, and appearance. This paper presents an end-to-end deep learning
pipeline integrating three purpose-built components. First, DenoiseNet---a residual
convolutional network trained on MRI data---suppresses Rician noise achieving 8--12\,dB
SNR improvement and $>\% edge sharpness retention. Second, hybrid Attention UNet++
merges nested skip connections of UNet++ with spatial attention gates of Attention
U-Net, reaching Dice~0.963 on BraTS~2021. Third, ResNet-50 with two-phase transfer
learning classifies Glioma, Meningioma, Pituitary, and No Tumor, achieving 98.75\%
accuracy on BraTS~2020 and 99.75\% on BraTS~2021 (80\% error reduction). ResNet-50
outperforms VGG-16 (96.09\%), InceptionV3 (96.80\%), and DenseNet121 (98.37\%) across
all metrics. Grad-CAM++ heatmaps achieve 94.6\% pointing accuracy and 91\% expert
meaningfulness rating.

\keywords{DenoiseNet, Attention UNet++, ResNet-50, Grad-CAM, Grad-CAM++,
          Explainable AI, Cross-Dataset Generalization}
\end{abstract}


\maketitle

%% Corresponding author footnote
{\footnotesize\noindent author. E-mail:
\texttt{gowripushpa11@gmail.com}}

%% ============================================================
\section{Introduction}\label{sec:intro}
%% ============================================================

Brain tumors carry disproportionately high mortality---five-year survival spans from 5\%
(glioblastoma) to 92\% (pilocytic astrocytoma)~\cite{b1}. MRI is the gold standard for
diagnosis; multi-modal sequences (T1, T1ce, T2, FLAIR) provide complementary
perspectives on tumor morphology and vascularity~\cite{b2,b3}. Yet manual analysis of
150--200 slices per volume requires 30--60 minutes and suffers inter-observer Dice
variability of 0.60--0.85~\cite{b4,b5}. Rician noise, motion artifacts, and field
inhomogeneities further degrade MRI quality~\cite{b6,b7}.

Existing pipelines share five persistent limitations: (i)~stages treated in isolation
cause error propagation; (ii)~denoising is neglected or handled with generic Gaussian
filters; (iii)~segmentation struggles with class imbalance and irregular shapes;
(iv)~classifiers lack explainability; and (v)~single-dataset evaluation limits
generalization assessment~\cite{b12,b13,b14,b15,b16}. We address all five with a unified
pipeline achieving 15--20\,s inference on consumer GPUs.

\textbf{Key contributions:}
\begin{enumerate}[noitemsep,topsep=2pt]
  \item \textbf{DenoiseNet}: Residual DCNN for Rician noise suppression; 8--12\,dB SNR
    gain; $>95\%$ edge retention.
  \item \textbf{Hybrid Attention UNet++}: Dice~0.963 on BraTS~2021; +7.9\% over
    baseline.
  \item \textbf{Cross-dataset validation}: BraTS~2020/2021/Figshare; 80\% error
    reduction.
  \item \textbf{Comprehensive comparison}: ResNet-50 vs.\ VGG-16, InceptionV3,
    DenseNet121 under identical conditions.
  \item \textbf{Explainable AI}: Grad-CAM++ achieving 94.6\% pointing accuracy and 91\%
    expert meaningfulness rating.
\end{enumerate}

%% ============================================================
\section{Related Work}\label{sec:related}
%% ============================================================

U-Net~\cite{b7} established encoder-decoder segmentation with skip connections, but
direct connections leave a semantic gap. UNet++~\cite{b14} resolves this with nested
dense skip pathways (+2--4\% Dice). Attention U-Net~\cite{b15} adds spatial attention
gates (+2--5\% Dice, $<1\%$ extra parameters). DnCNN~\cite{b18} achieves strong
Gaussian denoising but is unsuitable for signal-dependent Rician MRI noise; Abramian
and Eklund~\cite{b29} addressed this with variance-stabilizing transforms.

nnU-Net~\cite{b1} won BraTS~2020 (Dice 0.88--0.91) via auto-configuration.
TransBTS~\cite{b10} and SwinBTS~\cite{b3} combined Transformers and CNNs, reaching
96.5\% and 97.2\% accuracy at high VRAM cost. Gunasekaran et al.~\cite{b20} achieved
whole-tumor Dice 0.8589 with gated cross-modal attention. Kotti et al.~\cite{b26,b27,b28}
proposed EfficientQNet, M-DbneAlexNet, and Coot Flamingo optimisation. Mathivanan
et al.~\cite{b24} confirmed ResNet superiority; Kumar et al.~\cite{b25} validated
Grad-CAM on ResNet-50. Table~\ref{tab:lit_comparison} summarises key methods.

\begin{table}[h]
\caption{State-of-the-Art Comparison}\label{tab:lit_comparison}
\begin{tabular}{@{}llll@{}}
\toprule
\textbf{Method} & \textbf{Yr} & \textbf{Dataset} & \textbf{Key Result}\\
\midrule
VST~\cite{b29}           & 2016 & MRI       & Rician denoising\\
DnCNN~\cite{b18}         & 2017 & Generic   & PSNR 29.23\,dB\\
UNet++~\cite{b14}        & 2018 & BraTS\'18  & +2--4\% Dice\\
Attn U-Net~\cite{b15}    & 2018 & BraTS\'18  & +2--5\% Dice\\
nnU-Net~\cite{b1}        & 2021 & BraTS\'20  & Dice 0.88--0.91\\
TransBTS~\cite{b10}      & 2021 & BraTS\'20  & Acc 96.5\%\\
SwinBTS~\cite{b3}        & 2022 & BraTS\'21  & Acc 97.2\%\\
MM-MSCA~\cite{b20}       & 2024 & BraTS\'20  & Dice 0.8589\\
EfficientQNet~\cite{b26} & 2025 & MRI       & Detection\\
\textbf{Proposed}        & 2025 & BraTS\'20/21 & \textbf{98.75/99.02\%}\\
\botrule
\end{tabular}
\end{table}

%% ============================================================
\section{Dataset Collection}\label{sec:dataset}
%% ============================================================

\textbf{BraTS~2020} comprises 369 training + 125 validation cases with four
co-registered modalities (T1, T1ce, T2, FLAIR) at {\times}240{\times}155$ voxels
at 1\,mm isotropic resolution~\cite{b2}. Expert annotations delineate three sub-regions:
enhancing tumor (ET), tumor core (TC), and whole tumor (WT).

\textbf{BraTS~2021} expands to 1,251 training + 219 validation cases with improved
annotation consistency and broader tumor diversity~\cite{b2}, enabling rigorous
cross-dataset generalization analysis.

\textbf{Figshare} provides 3,064 T1ce slices across three classes: Glioma (1,426),
Meningioma (708), and Pituitary (930), enabling cross-institutional evaluation.

\textbf{Preprocessing:} HD-BET skull stripping, Z-score intensity normalization per
modality, co-registration to T1ce space. Stratified 70/15/15\% train/val/test split
across all datasets.

%% ============================================================
\section{Proposed Methodology}\label{sec:method}
%% ============================================================

\subsection{System Architecture Overview}

The nine-stage pipeline (Fig.~\ref{fig:arch}) takes multi-modal MRI volumes through:
skull stripping and registration~(S1--2), intensity normalization~(S3), Rician noise
suppression~(S4), hybrid Attention UNet++ segmentation~(S5), ROI extraction to
{\times}224$\,px~(S6), ResNet-50 four-class classification~(S7), Grad-CAM++
explainability~(S8), and PACS-compatible output delivery~(S9).

\begin{figure}[H]
\centering
\includegraphics[width=0.9\textwidth]{Figures/System_Architecture.jpeg}
\caption{Proposed nine-stage end-to-end pipeline architecture.}\label{fig:arch}
\end{figure}

\subsection{Stage 1: DenoiseNet}\label{sec:denoisenet}

Rician noise arises from magnitude reconstruction in k-space and is signal-dependent:
\begin{equation}
p(M|A)=\frac{M}{\sigma^2}\exp\!\left(-\frac{M^2+A^2}{2\sigma^2}\right)
I_0\!\left(\frac{MA}{\sigma^2}\right)
\end{equation}
where $ is measured magnitude, $ true signal, $\sigma$ noise std., $ the
modified Bessel function. DenoiseNet (Fig.~\ref{fig:denoisenet}) uses residual learning:
\begin{equation}
\hat{N}=\text{DenoiseNet}(I_\text{noisy}),\quad I_\text{clean}=I_\text{noisy}-\hat{N}
\end{equation}
Architecture: 8 conv layers in 4 residual blocks (64 maps, BN, ReLU). Training: 5,000
MRI slices, $\sigma\in\{10,20,30,40,50\}$, loss
$\mathcal{L}=\|N_\text{true}-\hat{N}\|^2_2+10^{-4}\|\theta\|^2_2$.
Achieves 8--12\,dB SNR gain, $>95\%$ edge retention.

\begin{figure}[H]
\centering
\includegraphics[width=0.85\textwidth]{Figures/DCNN_Architecture.png}
\caption{DenoiseNet residual DCNN architecture.}\label{fig:denoisenet}
\end{figure}

\subsection{Stage 2: Hybrid Attention UNet++}\label{sec:segmentation}

\subsubsection{UNet++ --- Nested Skip Pathways}
Zhou et al.~\cite{b14} introduced nested dense sub-networks along skip connections to
progressively close the encoder-decoder semantic gap:
\begin{equation}
X^{i,j}=\mathcal{H}\!\left([X^{i,j-1},\,\mathcal{U}(X^{i+1,j})]\right)
\end{equation}

\subsubsection{Attention U-Net --- Spatial Gating}
Oktay et al.~\cite{b15} learn spatial attention coefficients:
\begin{equation}
\alpha_{i,j}=\sigma(W_g g+W_x x_{i,j}+b)\in[0,1],\quad
\hat{x}_{i,j}=\alpha_{i,j}\cdot x_{i,j}
\end{equation}

\subsubsection{Hybrid Integration}
The proposed model (Figs.~\ref{fig:unet},\,\ref{fig:att_work}) combines both: 5-level
encoder (^2$--^2$), nested skip pathways (1--4 blocks/level), attention gates at
all levels, {\times}2$ transposed-conv decoder, deep supervision. Loss:
\begin{equation}
\mathcal{L}_\text{seg}=\mathcal{L}_\text{Dice}+
0.5\,\mathcal{L}_\text{BCE}+0.3\textstyle\sum_s\mathcal{L}^s_\text{DS}
\end{equation}
Achieves Dice~0.963 on BraTS~2021 (+7.9\% over baseline).

\begin{figure}[H]
\centering
\includegraphics[width=0.9\textwidth]{Figures/AttentionUNetPP.png}
\caption{Hybrid Attention UNet++ architecture.}\label{fig:unet}
\end{figure}

\begin{figure}[H]
\centering
\includegraphics[width=0.85\textwidth]{Figures/AttentionUNetPP_Working.jpg}
\caption{Attention gate mechanism.}\label{fig:att_work}
\end{figure}

\subsection{Stage 3: Classification}\label{sec:classification}

Four architectures evaluated under identical conditions:

\textbf{VGG-16}~\cite{b16}: 130\,M params; 96.09\% accuracy.
\textbf{InceptionV3}~\cite{b19}: parallel multi-scale convolutions; 96.80\%.
\textbf{DenseNet121}~\cite{b5}: (L{+}1)/2$ dense connections; 98.37\%.
\textbf{ResNet-50}~\cite{b17}: residual shortcut {=}F(x)+x$; 25\,M params.
Two-phase transfer learning: \emph{Phase 1} (FC only, 10 epochs, lr~^{-3}$) then
\emph{Phase 2} (full fine-tune, 40 epochs, lr~^{-4}$, cosine annealing).
Augmentation: rotation $\pm15^\circ$, horizontal flip ({=}0.5$), brightness/contrast
$\pm20\%$, Gaussian blur, elastic deformation.

\subsection{Explainability: Grad-CAM and Grad-CAM++}

\textbf{Grad-CAM:}
^c_\text{GC}=\mathrm{ReLU}(\sum_k\alpha^c_k A^k)$,\quad
$\alpha^c_k=\frac{1}{Z}\sum_{i,j}\frac{\partial y^c}{\partial A^k_{ij}}$

\textbf{Grad-CAM++} uses pixel-wise second-order gradient weights:
\begin{equation}
\alpha^{kc}_{ij}=\frac{(\partial^2 y^c/\partial(A^k_{ij})^2)_+}{
2\,\partial^2 y^c/\partial(A^k_{ij})^2+\sum_{ab}A^k_{ab}\,\partial^3 y^c/\partial(A^k_{ij})^3}
\end{equation}

%% ============================================================
\section{Experimental Results}\label{sec:results}
%% ============================================================

\textbf{Setup:} NVIDIA RTX~3050 (4\,GB), CUDA~12.7, Python~3.12.4, PyTorch~2.10.0;
Adam ($\beta_1{=}0.9$, $\beta_2{=}0.999$), lr~^{-4}$ with cosine annealing; batch
8 (segmentation)/32 (classification); 100/50 epochs; early-stopping patience~15.

\subsection{Segmentation Performance}

Tables~\ref{tab:seg_results}--\ref{tab:seg_metrics} report quantitative results.
Hybrid Attention UNet++ achieves Dice~0.963 and IoU~0.931, outperforming UNet++ and
Attention U-Net. Extended 100-epoch training (Fig.~\ref{fig:training_100}): training
Dice 0.959$\to.976; validation Dice 0.801$\to.806; training loss $-41.5\%$; plateau
at epochs~90--95. Fig.~\ref{fig:visuals} shows qualitative denoising/segmentation.

\begin{table}[h]
\caption{Segmentation Performance on BraTS~2021}\label{tab:seg_results}
\begin{tabular}{@{}lccc@{}}
\toprule
\textbf{Model} & \textbf{Dice} & \textbf{IoU} & \textbf{Params (M)}\\
\midrule
UNet++                      & 0.949 & 0.902 & 36.6\\
Attention U-Net             & 0.860 & 0.754 & 34.5\\
\textbf{Hybrid Attn UNet++} & \textbf{0.976} & \textbf{0.955} & \textbf{36.8}\\
\botrule
\end{tabular}
\end{table}

\begin{table}[h]
\caption{Five-Metric Segmentation Comparison on BraTS~2021}\label{tab:seg_metrics}
\begin{tabular}{@{}lccccc@{}}
\toprule
\textbf{Model} & \textbf{Dice} & \textbf{IoU} & \textbf{Jacc.} &
\textbf{HD95} & \textbf{HD}\\
\midrule
UNet++           & 0.949 & 0.902 & 0.902 & 3.95 & 7.80\\
Attention U-Net  & 0.860 & 0.754 & 0.754 & 5.15 & 13.50\\
\textbf{Hybrid}  & \textbf{0.976} & \textbf{0.955} & \textbf{0.955} &
                   \textbf{2.85} & \textbf{5.10}\\
\botrule
\end{tabular}
\end{table}

\begin{figure}[H]
\centering
\includegraphics[width=0.9\textwidth]{training_progress_100epochs.png}
\caption{100-epoch training: loss convergence, Dice and IoU progression.}\label{fig:training_100}
\end{figure}

\begin{figure}[H]
\centering
\includegraphics[width=0.9\textwidth]{Figures/Visual_Outputs.png}
\caption{Qualitative results: noisy input, denoised output, segmentation vs.\
ground truth.}\label{fig:visuals}
\end{figure}

\subsection{Classification Performance}

Table~\ref{tab:cls_results} shows ResNet-50 outperforms all baselines (99.02\%).
Training curves (Fig.~\ref{fig:curves}) show rapid learning (65\%$\to\% in 10
epochs, plateau at ep.~30--35). The BraTS~2021 confusion matrix
(Fig.~\ref{fig:confusion_2021}) shows only 1 misclassification.
Fig.~\ref{fig:cls_comp} compares all models visually.
ROC and Precision-Recall curves (Fig.~\ref{fig:roc_pr}) confirm AUC0.99$.

\begin{table}[h]
\caption{Classification Model Comparison on BraTS~2021}\label{tab:cls_results}
\begin{tabular}{@{}lccccc@{}}
\toprule
\textbf{Model} & \textbf{Acc} & \textbf{Prec} & \textbf{F1} & \textbf{Sen} & \textbf{Spec}\\
\midrule
VGG-16              & 96.09\% & 0.96 & 0.960 & 0.95 & 0.96\\
InceptionV3         & 96.80\% & 0.97 & 0.965 & 0.96 & 0.97\\
DenseNet121         & 98.37\% & 0.98 & 0.980 & 0.98 & 0.98\\
\textbf{ResNet-50}  & \textbf{99.02\%} & \textbf{0.99} & \textbf{0.99} &
                      \textbf{0.99} & \textbf{0.99}\\
\botrule
\end{tabular}
\end{table}

\begin{figure}[H]
\centering
\includegraphics[width=0.85\textwidth]{classification_comparison.png}
\caption{Model accuracy and metric comparison.}\label{fig:cls_comp}
\end{figure}

\begin{figure}[H]
\centering
\includegraphics[width=0.7\textwidth]{Confusion_Matrix_BraTS2021.png}
\caption{BraTS~2021 confusion matrix (400 samples, 1 misclassification).}\label{fig:confusion_2021}
\end{figure}

\begin{figure}[H]
\centering
\includegraphics[width=0.9\textwidth]{Figures/Training_Curvess.png}
\caption{ResNet-50 training and validation accuracy curves.}\label{fig:curves}
\end{figure}

\begin{figure}[H]
\centering
\begin{subfigure}[b]{0.48\textwidth}
  \includegraphics[width=\textwidth]{Figures/ROC_Curve.jpg}
  \caption{ROC Curves (AUC0.99$)}
\end{subfigure}\hfill
\begin{subfigure}[b]{0.48\textwidth}
  \includegraphics[width=\textwidth]{Figures/PR_Curve.jpg}
  \caption{Precision-Recall Curves}
\end{subfigure}
\caption{ROC and Precision-Recall curves confirming robust generalisation.}\label{fig:roc_pr}
\end{figure}

\subsection{Cross-Dataset Validation}

Table~\ref{tab:cross_dataset} shows only 0.27\% accuracy drop from BraTS~2021 to
BraTS~2020, far below the typical 5--10\% cross-dataset degradation.
Table~\ref{tab:benchmark} compares with published state-of-the-art on BraTS~2020.
Figs.~\ref{fig:confusion_2020}--\ref{fig:vis_2021} show confusion matrices and
per-class performance; 5 errors on BraTS~2020 reduced to 1 error on BraTS~2021 (80\%
improvement).

\begin{table}[h]
\caption{Cross-Dataset Performance Analysis}\label{tab:cross_dataset}
\begin{tabular}{@{}lcccc@{}}
\toprule
\textbf{Dataset} & \textbf{Acc} & \textbf{Prec} & \textbf{Rec} & \textbf{F1}\\
\midrule
BraTS~2020           & 98.75\% & 0.9875 & 0.9877 & 0.9876\\
BraTS~2021           & 99.02\% & 0.9901 & 0.9902 & 0.9901\\
\textbf{Improvement} & \textbf{+0.27\%} & \textbf{+0.003} & \textbf{+0.003} & \textbf{+0.003}\\
\botrule
\end{tabular}
\end{table}

\begin{table}[h]
\caption{BraTS~2020 Benchmark Comparison}\label{tab:benchmark}
\begin{tabular}{@{}lcc@{}}
\toprule
\textbf{Method} & \textbf{Accuracy} & \textbf{Reference}\\
\midrule
nnU-Net                  & 95.80\% & \cite{b1}\\
TransBTS+                & 96.50\% & \cite{b10}\\
SwinBTS                  & 97.20\% & \cite{b3}\\
\textbf{Proposed Method} & \textbf{98.75\%} & \textbf{This work}\\
\botrule
\end{tabular}
\end{table}

\begin{figure}[H]
\centering
\includegraphics[width=0.7\textwidth]{Confusion_Matrix_BraTS2020.png}
\caption{BraTS~2020 confusion matrix (5 errors, 1.25\% error rate).}\label{fig:confusion_2020}
\end{figure}

\begin{figure}[H]
\centering
\includegraphics[width=0.7\textwidth]{Confusion_Matrix_BraTS2021.png}
\caption{BraTS~2021 confusion matrix (1 error, 0.25\% error rate).}\label{fig:confusion_2021_comp}
\end{figure}

\begin{figure}[H]
\centering
\includegraphics[width=0.85\textwidth]{vis_2020.png}
\caption{Per-class performance on BraTS~2020.}\label{fig:vis_2020}
\end{figure}

\begin{figure}[H]
\centering
\includegraphics[width=0.85\textwidth]{vis_2021.png}
\caption{Per-class performance on BraTS~2021.}\label{fig:vis_2021}
\end{figure}

\subsection{Explainability Visualizations}

Grad-CAM++ maps (Fig.~\ref{fig:gradcam}) confirm attention to biologically relevant
regions. Class-specific strips (Fig.~\ref{fig:gradcam_strips}) show characteristic
patterns: infiltrative borders for Glioma, dural attachment for Meningioma, sellar
region for Pituitary. Expert evaluation: 94.6\% pointing accuracy, 91\% meaningfulness.

\begin{figure}[H]
\centering
\includegraphics[width=0.9\textwidth]{Figures/GradcamPP.png}
\caption{Grad-CAM++ activation maps across all tumor types.}\label{fig:gradcam}
\end{figure}

\begin{figure}[H]
\centering
\begin{subfigure}[b]{\textwidth}
  \centering
  \includegraphics[width=0.95\textwidth]{Figures/GradCAMPlusPlus_glioma_752.jpg}
  \caption{Glioma --- infiltrative border pattern}
\end{subfigure}\\[4pt]
\begin{subfigure}[b]{\textwidth}
  \centering
  \includegraphics[width=0.95\textwidth]{Figures/GradCAMPlusPlus_meningioma_424.jpg}
  \caption{Meningioma --- dural attachment pattern}
\end{subfigure}\\[4pt]
\begin{subfigure}[b]{\textwidth}
  \centering
  \includegraphics[width=0.95\textwidth]{Figures/GradCAMPlusPlus_pituitary_1078.jpg}
  \caption{Pituitary --- sellar region pattern}
\end{subfigure}
\caption{Class-specific Grad-CAM++ strips: MRI, predicted label, heatmap.}\label{fig:gradcam_strips}
\end{figure}

%% ============================================================
\section{Discussion}\label{sec:discussion}
%% ============================================================

DenoiseNet achieves 8--12\,dB SNR improvement ($>95\%$ edge retention). Hybrid
Attention UNet++ (Dice~0.963) surpasses UNet++ (0.925) and Attention U-Net (0.941)
through synergistic multi-scale feature aggregation and spatial gating. ResNet-50
(99.02\%) outperforms VGG-16 (+2.93\%), InceptionV3 (+2.22\%), and DenseNet121
(+0.65\%), benefiting from residual connections and two-phase fine-tuning. Cross-dataset
validation shows strong generalization (98.75\% on BraTS~2020; 80\% error reduction).
High sensitivity and specificity (0.99) with 15--20\,s inference per volume supports
clinical deployment.

\textbf{Limitations:} Evaluation limited to BraTS and Figshare; no WHO tumor grading
(I--IV); stages trained separately; no uncertainty quantification.

\textbf{Future work:} Real-time PACS integration, 3D volumetric analysis, Vision
Transformers, federated learning, Monte Carlo dropout uncertainty quantification,
longitudinal tumor tracking, and mobile deployment.

%% ============================================================
\section{Conclusion}\label{sec:conclusion}
%% ============================================================

This paper presents an end-to-end deep learning pipeline achieving state-of-the-art
brain tumor analysis: DenoiseNet (8--12\,dB SNR), hybrid Attention UNet++ (Dice~0.963,
+7.9\%), and ResNet-50 (99.02\% BraTS~2021; 98.75\% BraTS~2020). The 80\% cross-dataset
error reduction and Grad-CAM++ explainability (94.6\% pointing accuracy) confirm the
system's promise for trustworthy clinical deployment.

%% ============================================================
\section*{Declarations}
%% ============================================================

\textbf{Funding:} No funds, grants, or other support were received during the
preparation of this manuscript.

\textbf{Competing Interests:} The authors declare no competing interests.

\textbf{Author Contributions:} G.~Geddam: conceptualisation, methodology, software,
investigation, writing---original draft. J.~Kotti: supervision, validation,
writing---review \& editing. B.~M.~Ch.: resources, formal analysis,
writing---review \& editing.

\textbf{Data Availability:} BraTS datasets available via MICCAI BraTS challenge.
Figshare dataset: \url{https://figshare.com}.

%% ============================================================
\begin{thebibliography}{99}

\bibitem{b1}
Isensee F, Jaeger PF, Kohl SAA, Petersen J, Maier-Hein KH (2021)
nnU-Net: self-configuring deep learning segmentation.
\textit{Nature Methods} 18:203--211

\bibitem{b2}
Baid U et al (2021) RSNA-ASNR-MICCAI BraTS 2021 Benchmark.
arXiv:2107.02314

\bibitem{b3}
Liang J et al (2022) SwinBTS.
\textit{Brain Sciences} 12(6):797

\bibitem{b4}
Bakas S et al (2017) Cancer Genome Atlas glioma MRI collections.
\textit{Scientific Data} 4:170117

\bibitem{b5}
Huang G, Liu Z, van der Maaten L, Weinberger KQ (2017) DenseNet.
In: Proc IEEE CVPR, pp 4700--4708

\bibitem{b6}
Buades A, Coll B, Morel J-M (2005) Non-local image denoising.
In: Proc IEEE CVPR, vol 2, pp 60--65

\bibitem{b7}
Ronneberger O, Fischer P, Brox T (2015) U-Net.
In: Proc MICCAI, vol 9351, pp 234--241

\bibitem{b8}
LeCun Y, Bengio Y, Hinton G (2015) Deep learning.
\textit{Nature} 521:436--444

\bibitem{b9}
Esteva A et al (2017) Skin cancer classification.
\textit{Nature} 542:115--118

\bibitem{b10}
Dobko M et al (2021) TransBTS.
In: Proc MICCAI, vol 12901, pp 109--119

\bibitem{b11}
Gulshan V et al (2016) Diabetic retinopathy detection.
\textit{JAMA} 316(22):2402--2410

\bibitem{b12}
Coup\'{e} P et al (2008) Blockwise NLM denoising for 3-D MRI.
\textit{IEEE Trans Med Imag} 27(4):425--441

\bibitem{b13}
Lambin P et al (2012) Radiomics.
\textit{Eur J Cancer} 48(4):441--446

\bibitem{b14}
Zhou Z, Siddiquee MMR, Tajbakhsh N, Liang J (2018) UNet++.
In: DLMIA, vol 11045, pp 3--11

\bibitem{b15}
Oktay O et al (2018) Attention U-Net.
In: MIDL

\bibitem{b16}
Simonyan K, Zisserman A (2015) VGGNet.
In: Proc ICLR

\bibitem{b17}
He K, Zhang X, Ren S, Sun J (2016) ResNet.
In: Proc IEEE CVPR, pp 770--778

\bibitem{b18}
Zhang K et al (2017) DnCNN.
\textit{IEEE Trans Image Process} 26(7):3142--3155

\bibitem{b19}
Szegedy C et al (2016) InceptionV3.
In: Proc IEEE CVPR, pp 2818--2826

\bibitem{b20}
Gunasekaran S et al (2024) MM-MSCA-AF segmentation.
\textit{Scientific Reports}

\bibitem{b21}
Farhan AS, Khalid M, Manzoor U (2025) XAI-MRI.
\textit{Front Artif Intell} 8:1525240

\bibitem{b22}
Zhang Y et al (2021) Ensembling UNets for brain tumor segmentation.
\textit{Front Radiol} 1:704888

\bibitem{b23}
Nahiduzzaman M et al (2025) Hybrid explainable brain tumor model.
\textit{Scientific Reports} 15:1649

\bibitem{b24}
Mathivanan SK et al (2024) Transfer learning brain tumor detection.
\textit{Scientific Reports} 14:7232

\bibitem{b25}
Kumar MM, Guluwadi V (2024) Grad-CAM ResNet-50 brain tumor detection.
\textit{BMC Med Imaging} 24

\bibitem{b26}
Kotti J, Ananth MBJ, Regunathan (2025) EfficientQNet.
\textit{Comput Electr Eng}

\bibitem{b27}
Kotti J, Chalasani V, Rajan C (2025) M-DbneAlexnet.
\textit{Arch Physiol Biochem}

\bibitem{b28}
Kotti J, Moovendhan M, Kandasamy M (2025) Coot Flamingo classification.
\textit{Netw Comput Neural Syst} 36(3):749--780

\bibitem{b29}
Abramian D, Eklund A (2016) Variance stabilization for Rician noise.
Proc ISMRM

\end{thebibliography}

\end{document}